\begin{abcliste}
\abc
Der Hebelarm ist der Abstand von D zur Wirkungslinie. Alle vier Kräfte sind gleich gross und greifen im selben Punkt an, es entscheidet also allein der Hebelarm:
\begin{itemize}
 \item Die Wirkungslinie von $F_1$ geht durch D: Hebelarm null.
 \item $F_2$ steht senkrecht auf DP: Der Hebelarm ist die ganze Strecke DP, so lang wie er überhaupt sein kann.
 \item $F_3$ und $F_4$ sind gleich stark gegen DP geneigt: Ihre Wirkungslinien haben denselben Abstand von D, kleiner als DP und grösser als null.
\end{itemize}
\[
 \result{M_1 < M_3 = M_4 < M_2}
\]

\abc
$F_1$ dreht die Platte gar nicht. $F_2$ und $F_3$ ziehen P nach oben und drehen die Platte \result{\text{im Gegenuhrzeigersinn}}, $F_4$ zieht P nach unten und dreht sie \result{\text{im Uhrzeigersinn}}. $M_3$ und $M_4$ sind gleich gross, haben aber entgegengesetzten Drehsinn.

\abc
Sei $r$ der Abstand DP und $\alpha$ der Winkel zwischen Kraft und DP.
\begin{align}
 M_1 &= \result{0} \\
 M_2 &= \MzwF = \FP\cdot\rP = \Mzw \approx \result{\MzwP{0}{2}}
\end{align}
Für $F_3$ und $F_4$ ist der Hebelarm
\begin{align}
 d &= \dschF = \rP\cdot\sin\al = \dsch \approx \dschP{-2}{2}
\end{align}
und damit
\begin{align}
 M_3 = M_4 &= \MschF = \FP\cdot\rP\cdot\sin\al = \Msch \approx \result{\MschP{0}{2}}
\end{align}
Man kann auch die Kraft zerlegen: Nur die Komponente senkrecht zu DP, $F\sin\alpha$, hat einen Hebelarm; die Komponente entlang DP geht durch D.
\end{abcliste}
